% ADD scatterplot matrices, and yes you can do them in Python

\chapter{Spyder, review, misc tricks}

\section{Spyder}
We'll be using Spyder this semester: a Python-based data analysis environment
written especially for Data Science and scientific programming.

It comes with the Anaconda distribution, as does pretty much everything else
we'll use this semester (scikit-learn, NumPy, SciPy, Pandas, Matplotlib,
\textit{etc.}) Download it from here: \url{https://www.anaconda.com/download}. 

\textbf{Note:} make sure to get Python \textbf{3.*}, not \textbf{2.*}!! (They
are \textit{not} mutually compatible.) The number(s) after the dot don't
matter so much. But before the dot \textit{must} be a \textbf{``3''}.

Don't worry: all the stuff you learned using Python Jupyter Notebooks still
applies! The only difference is that using an IDE like Spyder is not
Web-based: it's entirely offline, responsive, and more powerful. You will be
creating your own \texttt{.py} files directly.


\section{Review (?) from DATA 101}

\begin{compactitem}
\itemsep.1em
\item Data stuff
    \begin{compactitem}
    \item Contingency tables
    \item Weighted averages (and averages of averages)
    \item Scatterplots
    \item Histograms
    \item Boxplots
    \item Quantiles
    \end{compactitem}

\item Python stuff
    \begin{compactitem}
    \item Lists and dictionaries
    \item \texttt{for} loops
    \item Modifying in-place vs. returning a copy
    \item Pandas Series \& DataFrames
        \begin{compactitem}
        \itemsep.1em
        \item The ``index''
        \item \texttt{.iloc} vs. \texttt{.loc} vs. neither
        \item Single ints/labels vs. slices vs. lists
        \item Queries
        \end{compactitem}
    \end{compactitem}
\end{compactitem}

\section{Some more advanced extensions to the above}

\subsection{Multi-list iteration}

Suppose you have:
%\begin{Verbatim}[fontsize=\small,samepage=true,frame=single]
    episode7 = [ 9, 6, 10, 10, 9, 10, 10, 10 ]
    episode8 = [ 9, 7,  4,  8, 8, 10, 10,  7 ]
%\end{Verbatim}
and you want a new list \texttt{avgSW} that gives the average rating each
viewer gave the two movies.

\begin{enumerate}
\itemsep.1em
\item One way to do it (caveman style):
\begin{Verbatim}[fontsize=\small,samepage=true,frame=single]
avgSW = []
for i in range(len(episode7)):
    avgSW.append((episode7[i]+episode8[i])/2)
\end{Verbatim}

\item Another way to do it (modern idiomatic Python):
\begin{Verbatim}[fontsize=\small,samepage=true,frame=single]
avgSW = []
for e7, e8 in zip(episode7, episode8):
    avgSW.append((e7+e8)/2)
\end{Verbatim}

\end{enumerate}


\subsection{Pandas DF updating}

Sometimes you can run into trouble when trying to update a Pandas DataFrame.
For instance, take this DataFrame:

\begin{Verbatim}[fontsize=\small,samepage=true,frame=single]
>>> print(characters)
                           rank  saber   side
    name                                     
    Darth Vader       Sith Lord    red   Dark
    Obi-wan Kenobi  Jedi Master   blue  Light
    Yoda            Jedi Master  green  Light
    Luke Skywalker         Jedi  green  Light
    Kylo Ren         Apprentice    red   Dark
    Rey              Apprentice   blue  Light
\end{Verbatim}

This query works just fine:

\begin{Verbatim}[fontsize=\small,samepage=true,frame=single]
>>> characters.loc['Kylo Ren':'Rey']['rank']
    name
    Kylo Ren    Apprentice
    Rey         Apprentice
    Name: rank, dtype: object
\end{Verbatim}

so naturally you'd think you could do this:
\begin{Verbatim}[fontsize=\small,samepage=true,frame=single]
>>> characters.loc['Kylo Ren':'Rey']['rank'] = 'Jedi'
\end{Verbatim}

Alas, no: \texttt{A value is trying to be set on a copy of a slice from a
DataFrame.}

The solution is to use this syntax:

\begin{Verbatim}[fontsize=\small,samepage=true,frame=single]
>>> characters.loc['Kylo Ren':'Rey','rank'] = 'Jedi'
\end{Verbatim}

(Your mileage may vary.)

\subsection{Contingency tables in Pandas}

\subsubsection{\texttt{.value\_counts()}}
You can do a quick and dirty contingency table in Pandas using
\texttt{.value\_counts()} (possibly in conjunction with \texttt{.groupby()}):

\begin{Verbatim}[fontsize=\small,samepage=true,frame=single]
>>> characters.side.value_counts()
    Light    4
    Dark     2
    Name: side, dtype: int64

>>> characters.groupby('rank').side.value_counts()
    rank         side 
    Apprentice   Dark     1
                 Light    1
    Jedi         Light    1
    Jedi Master  Light    2
    Sith Lord    Dark     1
    Name: side, dtype: int64
\end{Verbatim}

\subsubsection{\texttt{pd.crosstab()}}
The \texttt{pd.crosstab()} function is a bit easier on the eyes, though, and
can compute margins:

\begin{Verbatim}[fontsize=\small,samepage=true,frame=single]
>>> side_saber = pd.crosstab(characters.side, columns=characters.saber, 
                                                             margins=True)
    saber  blue  green  red  All
    side                        
    Dark      0      0    2    2
    Light     2      2    0    4
    All       2      2    2    6
\end{Verbatim}

Get overall percentages with:
\begin{Verbatim}[fontsize=\small,samepage=true,frame=single]
>>> side_saber / side_saber.loc['All']['All'] * 100
    saber       blue      green        red         All
    side                                              
    Dark    0.000000   0.000000  33.333333   33.333333
    Light  33.333333  33.333333   0.000000   66.666667
    All    33.333333  33.333333  33.333333  100.000000
\end{Verbatim}

or by column by doing:
\begin{Verbatim}[fontsize=\small,samepage=true,frame=single]
>>> side_saber / side_saber.loc['All'] * 100
    saber   blue  green    red         All
    side                                  
    Dark     0.0    0.0  100.0   33.333333
    Light  100.0  100.0    0.0   66.666667
    All    100.0  100.0  100.0  100.000000
\end{Verbatim}

By rows is a bit trickier:
\begin{Verbatim}[fontsize=\small,samepage=true,frame=single]
>>> side_saber.div(side_saber['All'],axis=0) * 100
    saber       blue      green         red    All
    side                                          
    Dark    0.000000   0.000000  100.000000  100.0
    Light  50.000000  50.000000    0.000000  100.0
    All    33.333333  33.333333   33.333333  100.0
\end{Verbatim}

% Significant?:
% import scipy.stats.chi2_contingency
% chi2_contingency(side_saber)

\vspace{4in}
\pagebreak
\subsection{Notched boxplots}

You can add the argument ``\texttt{notch=True}'' to
``\texttt{aDataFrame.boxplot()}'' to get a ``notched'' boxplot. This visually
indicates whether or not there is a stat sig diff between groups, based on
\textit{whether the notched areas overlap}:

\begin{Verbatim}[fontsize=\small,samepage=true,frame=single]
>>> students.boxplot('sat',by='gender',notch=True)
>>> students.boxplot('height',by='gender',notch=True)
\end{Verbatim}

\begin{center}
\includegraphics[width=0.45\textwidth]{sat.pdf}
\includegraphics[width=0.45\textwidth]{height.pdf}
\end{center}

